We formalize our problem as follows. 
An employer accesses an infinite pool of candidates 
(indexed by $i \in \mathbb{N}^+$), 
each of which has some (hidden) \textit{skill level} $y_i\in\{0,+1\}$,
which denote \emph{unskilled} and \emph{skilled}, respectively. 
Underlying worker skill levels $y_i$ are sampled independently
from a Bernoulli distribution with parameter $p$.
An employer can access information about the $i$-th candidate
through a sequence of $\tau$ tests, 
which are conditionally independent given $y_i$. 
Each \textit{test result}, $\hat{y}_{i,j}\in \{0,+1\}$ 
disagrees with the ground truth skill with probability
$\Pr[\hat{y}_{i,j} \neq y_i]=\frac{1-\sigma}{2}$, 
where $\sigma\in(0,1)$, i.e., $\hat{y}_{i,j}=y_i\xor Br(\frac{1-\sigma}{2})$\footnote{$\xor$ is the XOR operation between two binary random variables, and therefore $\hat{y}_{i,j}$ is also a random variable.}.
For convenience, we denote the noise level as $\eta=\frac{1-\sigma}{2}\in (0,\frac{1}{2})$.
We say that a test result $\hat{y}_{i,j}$ is \textit{flipped} 
if $\hat{y}_{i,j} \ne y_i$, 
and the number of flipped results for a given candidate 
is denoted by 
$Z_\tau^\eta$ is $Z_\tau^\eta=\sum_{j=1}^\tau \mathbb{I}(\hat{y}_{i,j}\ne y_i)$, 
where $\mathbb{I}(\cdot)$ is the indicator function. 


A \emph{selection criterion} is a mapping 
between test results to actions: $\text{\tt Select}(\hat{y}_{i,1},\dots,\hat{y}_{i,\tau_i}) \in \{0,1\}$, where $0$ means \emph{reject} 
and $1$ means \emph{accept} (hire). 
A \emph{policy} $\pi$ sets the selection criteria 
based on $\sigma, p$ and other possible constraints 
such as probability to hire, error probability, etc. 
A \emph{randomized threshold policy} is a policy $\pi$ with parameters $(\tau,\theta,r)$ such that $\pi(\hat{y}_{i,1},\dots,\hat{y}_{i,\tau_i})=1$ 
for $S_{\tau}:=\sum_{i=1}^{\tau} \hat{y}_{i,j} > \theta$,
$\pi(\hat{y}_{i,1},\dots,\hat{y}_{i,\tau_i})=0$ 
for $S_{\tau}< \theta$, and for $S_{\tau}= \theta$ the probability that $\pi(\hat{y}_{i,1},\dots,\hat{y}_{i,\tau_i})=1$ is $r$.
We call a policy $\pi$ a \emph{threshold policy} if $r=1$.
In a \textit{dynamic policy}, 
rather than setting a fixed number of tests per candidate,
the employer may decide after each test whether to 
\emph{accept}, \emph{reject}, 
or to perform an additional test, i.e., $\pi(\hat{y}_{i,1},\dots,\hat{y}_{i,\tau_i}) \in \{0,1,\texttt{more}\}$. 
Note that for a dynamic policy, the number of tests $\tau$
is a random variable determined based on the tests' outcomes.
When designing a policy, one must carefully consider 
the balance between the following desiderata:
\begin{enumerate}
\item \textbf{Minimize False Discovery Rate (FDR)}---the fraction of unskilled workers 
among the accepted candidates, i.e., $\text{FDR}:=\Pr[y_i=0|\pi(\hat{y}_{i,1},\dots,\hat{y}_{i,\tau})=+1]$.
\item \textbf{Minimize False Omission Rate (FOR)}---the fraction of skilled workers among the rejected candidates, i.e., $\text{FOR}:=\Pr[y_i=+1|\pi(\hat{y}_{i,1},\dots,\hat{y}_{i,\tau})=0]$. 
\item \textbf{Minimize False Negatives (FN)}---the amount of skilled workers that are classified as unskilled.
\item \textbf{Minimize False Positives (FP)}---the amount of unskilled workers that are classified as skilled.
\item \textbf{Ratio of accept probability and the number of tests}---the number 
of tests performed per candidate hired, 
using a parameter $B>1$, we have  $\frac{\tau}{B} \leq \Pr[\pi(\hat{y}_{i,1},\dots,\hat{y}_{i,\tau})=+1]$.
\end{enumerate}
For any fixed 
number of tests $\tau$, 
increasing the threshold $\theta$ of a threshold policy
decreases FDR while increasing FOR.

\noindent{\bf Loss:} 
To handle the trade-off between the false positives, 
(i.e., when an unskilled candidate is accepted) 
and false negatives 
(i.e., when a skilled candidate is rejected), 
we introduce an $\alpha$-loss, 
paramaterized by $\alpha \in [0,1]$ and defined as follows:
\begin{equation*}
\ell_{\alpha}(b_1,b_2)=\alpha \mathbb{I}[b_1=0,b_2=1] + (1-\alpha)\mathbb{I}[b_1=1,b_2=0]
\end{equation*}
where $\mathbb{I}[\cdot]$ is the indicator function and $b_1,b_2\in\{0,1\}$. The expected loss of a policy $\pi$ is,
\begin{equation}\label{loss}
    l_\alpha(\pi)=E[\ell_{\alpha}(y_i,\pi(\hat{y}_{i,1},\dots,\hat{y}_{i,\tau}))] 
\end{equation}
where the expectation is over the type of the candidates $y_i$, the test results $\hat{y}_{i,j}$, and the decisions of $\pi$.